\documentclass[11pt,a4paper]{../_shared/altacv}

\geometry{left=1.5cm,right=1.5cm,top=1.5cm,bottom=1.5cm,columnsep=1.5cm}
\usepackage[utf8]{inputenc}
\usepackage[T1]{fontenc}
\usepackage[french]{babel}
\usepackage{paracol}
\usepackage{xcolor}
\usepackage{hyperref}
\usepackage{fontawesome5}
\usepackage{setspace}

% Couleurs exactes du CV
\definecolor{SlateGrey}{HTML}{2E2E2E}
\definecolor{LightGrey}{HTML}{666666}
\definecolor{DarkPastelRed}{HTML}{8AC0D9}
\definecolor{PastelRed}{HTML}{00B0DA}
\definecolor{GoldenEarth}{HTML}{B4AD0C}
\colorlet{name}{black}
\colorlet{tagline}{PastelRed}
\colorlet{heading}{DarkPastelRed}
\colorlet{headingrule}{GoldenEarth}
\colorlet{subheading}{PastelRed}
\colorlet{accent}{PastelRed}
\colorlet{emphasis}{SlateGrey}
\colorlet{body}{LightGrey}

\renewcommand{\familydefault}{\sfdefault}

% En-tête identique au CV
\name{BOILEAU Guillaume}
%\tagline{R\&D Engineer and Data Project Manager}
%\photoL{2.8cm}{../_shared/moi}
\personalinfo{
  \email{guillaume.boileaupro@gmail.com}
  \phone{+33 6 59 94 85 84}
  \location{Alpes Maritimes, France}
  \linkedin{guillaume-boileau-phd-891b81265}
  \researchgate{Guillaume-Boileau-2}
  \orcid{0000-0002-3576-6968}
}

\setlength{\parskip}{0.75\baselineskip}
%\linespread{1.15}

\begin{document}
\makecvheader

\cvsection{Lettre de motivation}
{\color{accent}\textbf{Objet :}} \textit{Candidature au poste de Senior Data Engineer -- Réf. 402255}

{\color{body}

Madame, Monsieur,

Actuellement ingénieur logiciel et data chez VEV Platform Services France à Valbonne, je souhaite vous proposer ma candidature au poste de Senior Data Engineer chez Capgemini.

Titulaire d’un doctorat en physique, j’ai développé au cours de mon parcours une expertise solide dans le traitement de données complexes, la conception de pipelines robustes, la modélisation, l’analyse statistique et le développement logiciel. Mon expérience s’est construite à l’interface entre recherche scientifique, ingénierie des données et développement applicatif, dans des environnements exigeants où la fiabilité des traitements, la qualité des données et la capacité à transformer des besoins complexes en solutions opérationnelles étaient centrales.

Dans mes fonctions actuelles, je développe des services backend et des outils de traitement de données pour des infrastructures techniques connectées. J’interviens sur l’intégration de flux issus de systèmes variés, la structuration et l’exploitation de données opérationnelles, le développement d’API, ainsi que la mise en place de services robustes destinés à être exploités dans des environnements de production. Cette expérience m’a permis de renforcer ma pratique de l’industrialisation, de l’intégration de services, de la documentation technique et du dialogue avec des interlocuteurs aux besoins variés.

Auparavant, dans le cadre de mes activités de recherche et de R\&D, j’ai conçu et développé plusieurs pipelines scientifiques en Python pour l’analyse, la validation et la consolidation de jeux de données complexes à grande échelle. Ces travaux m’ont conduit à structurer des chaînes de traitement reproductibles, à modéliser des données hétérogènes, à développer des outils d’automatisation et à exploiter des environnements Linux, HPC et distribués. J’y ai acquis une forte culture de la qualité des traitements, de la traçabilité, de la robustesse méthodologique et de la performance, qui me paraît particulièrement pertinente pour des projets de data engineering à forte exigence technique.

Je me reconnais pleinement dans les missions proposées par Capgemini : conception et optimisation de pipelines de données, intégration de sources multiples, contribution à des plateformes data modernes, mise en qualité et documentation, ainsi qu’appui à des cas d’usage data, analytics et IA. Mon profil associe en effet une forte capacité analytique, une expérience concrète du développement logiciel et une réelle aisance à travailler sur des données complexes dans une logique de production et de valeur métier.

Mes compétences principales portent aujourd’hui sur Python, SQL, le développement de services et d’API, Linux, Docker, Git, l’analyse de données, la modélisation statistique et le machine learning. J’ai également une bonne familiarité avec les architectures logicielles modernes, les environnements conteneurisés et les logiques d’intégration de services. Mon parcours m’a aussi appris à monter rapidement en compétence sur de nouveaux écosystèmes techniques, ce qui me permettrait de contribuer efficacement à des environnements data plus avancés intégrant cloud, orchestration, gouvernance et plateformes de type Data Lake ou Lakehouse.

Au-delà des compétences techniques, j’accorde une grande importance au travail collectif, à la pédagogie et à la compréhension du besoin. Mes expériences m’ont amené à collaborer avec des profils très variés, à expliciter des sujets techniques complexes, à proposer des solutions pragmatiques et à accompagner des projets dans la durée. Cette capacité à faire le lien entre enjeu métier, exigence technique et qualité d’exécution me semble pleinement en phase avec l’approche attendue pour ce poste.

Rejoindre Capgemini représenterait pour moi l’opportunité de mettre cette double culture data et développement au service de projets ambitieux, innovants et à fort impact, dans un environnement stimulant où l’expertise technique, le cloud, l’IA et la transformation des usages occupent une place centrale.

Je vous remercie pour l’attention portée à ma candidature et me tiens naturellement à votre disposition pour tout échange complémentaire.

Veuillez agréer, Madame, Monsieur, l’expression de mes salutations distinguées.

\textbf{Guillaume Boileau}

}


\end{document}